\renewcommand{\theposttitle}{How expertise hurts your creativity and the ancient Zen Buddhist mindset that changes it}
\renewcommand{\thepostdate}{2023-12-19}
\vspace*{\fill}
\begin{center}
\includegraphics[width=\textwidth, height=0.55\textheight, keepaspectratio]{C:/Users/James/Documents/GitHub/VisualIdeasBook/data/images/c9b024bb932aaba345699a745ef657f6.png}
\end{center}
\vspace{2mm}
\begin{center}
\begin{minipage}{0.92\textwidth}
\small
\textbf{The definition of integrity. }

One of my favorite things to journal on is defining what a word means to me.

I started doing this after studying for the GRE (a test students have to take in the states to go to graduate school).

It made me realize there were lots of words I knew how to use in a sentence, but could only explain vaguely.

For example, take a simple word like “love”.

If you’ve ever read \textbf{The Five Love Languages}, then you know people have different (often subconscious) ways they define love.

And if you define what love means to you and see how it aligns with your partner’s definition, it’s pretty helpful to communicate more effectively in your relationship.

In the same vein, I’ve found clarifying how you understand other vague words (like integrity) helps align a deeper relationship with yourself.

I’m curious, is your definition of integrity the same or different?
\end{minipage}
\end{center}
\vspace*{\fill}